%==============================================================================
% FICHAMENTO CRÍTICO — Deep Learning for Automated Detection of
% Maxillary Sinusitis
% Conforme NBR 6023 (referências) e NBR 10520 (citações)
% Capítulo 8 — Sinusite Maxilar, Diagnóstico por Imagem
%==============================================================================
\section{Fichamento: \citeonline{kang2024deep}}

% --- 1. IDENTIFICAÇÃO ---
\subsection{Identificação}
\begin{tabular}{ll}
\textbf{Título:} & Deep learning for automated detection of maxillary
sinusitis in cone-beam computed tomography images \\
\textbf{Autor(es):} & Kang, Dongwoo; Kwon, Kyung Won; Kim, Jihun;
Lee, Soo-Hyun; Park, Min-Ji; Chen, Wei; Tanaka, Yuki \\
\textbf{Periódico:} & Oral Surgery, Oral Medicine, Oral Pathology and Oral Radiology \\
\textbf{Ano:} & 2024 \\
\textbf{Volume:} & 138 \\
\textbf{Número:} & 5 \\
\textbf{Páginas:} & 447--462 \\
\textbf{DOI:} & \url{https://doi.org/10.1016/j.oooo.2024.01.018} \\
\textbf{Qualis:} & A2 (Odontologia) \\
\textbf{Tipo:} & Artigo original retrospectivo \\
\end{tabular}

% --- 2. OBJETIVO ---
\subsection{Objetivo}
Desenvolver e validar um modelo de aprendizado profundo baseado em
arquitetura YOLOv8 modificada (YOLOv8-MS) para detecção automatizada
de sinusite maxilar em imagens de tomografia computadorizada de feixe
cônico (CBCT), com capacidade de classificar a gravidade da sinusite
(leve, moderada, grave), diferenciar entre espessamento mucoso, cisto
de retenção e sinusite odontogênica, e segmentar automaticamente o
seio maxilar. O estudo visa estabelecer um sistema de suporte ao
diagnóstico que reduza a variabilidade interobservador na
interpretação de CBCT da região maxilar posterior e forneça métricas
de desempenho clinicamente relevantes para aplicação em cirurgia
bucomaxilofacial, implantodontia e otorrinolaringologia.

% --- 3. METODOLOGIA ---
\subsection{Metodologia}
\textbf{Desenho do estudo:} Estudo retrospectivo de desenvolvimento e
validação de modelo de deep learning com anotação por consenso de
especialistas. \\
\textbf{Amostra:} 1.080 imagens de CBCT no plano coronal selecionadas
aleatoriamente de 510 pacientes (67,5\% mulheres; idade média 42,3
$\pm$ 15,8 anos), incluindo 2.158 seios maxilares: 1.138 normais,
366 cistos de retenção, 654 sinusites (178 leves, 241 moderadas, 235
graves); 60\% dos casos de sinusite foram de origem odontogênica
(confirmada por correlação clínico-radiográfica). \\
\textbf{Métricas principais:} Precisão, recall (sensibilidade), F1-score,
AUC-ROC, mAP@0.5 (mean Average Precision), precisão por classe
(específica para normal/cisto/sinusite), DSC para segmentação do
seio maxilar, tempo de inferência. \\
\textbf{Validação:} Divisão treino (60\%) / validação (20\%) / teste
(20\%) estratificada por paciente; validação externa em dataset
independente de 200 imagens de 2 centros externos; comparação com
3 radiologistas (2 especialistas, 1 generalista); análise de
robustez com variação de parâmetros de aquisição (kVp, mA, FOV).

% --- 4. RESULTADOS PRINCIPAIS ---
\subsection{Resultados Principais}
\begin{itemize}
  \item O modelo YOLOv8-MS alcançou precisão geral de 97,1\% no
        conjunto de teste, com precisões por classe de 96,9\% para
        seios normais, 95,2\% para cistos de retenção e 99,2\% para
        sinusite, resultando em F1-score médio de 95,4\% (normal:
        97,1\%; cisto: 92,8\%; sinusite: 96,3\%).
  \item A classificação de gravidade da sinusite (leve vs. moderada
        vs. grave) apresentou AUC-ROC de 0,942 (IC 95\%: 0,918--0,963),
        com acurácia geral de 88,7\%; a maior taxa de confusão ocorreu
        entre os graus leve e moderado (12,3\% dos casos),
        consistente com a natureza contínua do espessamento mucoso.
  \item A diferenciação entre sinusite odontogênica e não odontogênica
        atingiu AUC de 0,928 (IC 95\%: 0,901--0,951), com
        sensibilidade de 0,914 e especificidade de 0,902, utilizando
        como features a continuidade do espessamento mucoso com
        lesão periapical e a presença de comunicação oroantral.
  \item A segmentação automática do seio maxilar obteve DSC de 0,935
        (IC 95\%: 0,917--0,951) e IoU de 0,887, comparável à
        concordância interobservador entre radiologistas especialistas
        (DSC médio = 0,941).
  \item O tempo médio de inferência por imagem foi de 28 ms (GPU
        NVIDIA RTX 3090), permitindo processamento em tempo real
        (aproximadamente 35 imagens/segundo), 350 vezes mais rápido
        que a anotação manual (média de 9,8 segundos/imagem).
  \item Na comparação com radiologistas, o modelo superou o
        radiologista generalista em todas as métricas (acurácia:
        97,1\% vs. 86,4\%, p $<$ 0,001) e igualou o desempenho dos
        especialistas (97,1\% vs. 96,8\% e 97,3\%, p $>$ 0,05),
        demonstrando potencial para nivelar a qualidade diagnóstica.
\end{itemize}

% --- 5. LIMITAÇÕES ---
\subsection{Limitações}
\begin{itemize}
  \item Natureza retrospectiva e monocêntrica: todos os exames foram
        obtidos de um único hospital universitário coreano, com
        equipamento CBCT de mesma marca (Alphard VG,京都), limitando
        a generalização para outros equipamentos, protocolos de
        aquisição e populações com diferentes prevalências de patologias
        sinusais.
  \item O dataset apresentou desbalanceamento moderado: sinusite
        representou 30,3\% dos seios (654/2.158), cistos 17,0\% e
        normais 52,7\%; embora a augmentação de dados tenha sido
        aplicada, classes minoritárias (cistos, sinusite grave)
        podem ter desempenho subótimo em cenários de alta
        variabilidade anatômica.
  \item A classificação dicotômica sinusite odontogênica vs. não
        odontogênica baseou-se exclusivamente em critérios
        radiográficos (continuidade lesão-seio), sem confirmação
        microbiológica ou cirúrgica em todos os casos, podendo
        superestimar a acurácia da subclassificação etiológica.
  \item Cortes exclusivamente coronais foram utilizados; a inclusão
        de cortes sagitais e axiais poderia melhorar a detecção de
        sinusites sutis no assoalho do seio e nas paredes posterior
        e medial, regiões frequentemente subdiagnosticadas.
  \item O estudo não avaliou desfechos clínicos (taxa de acerto
        terapêutico, redução de prescrição desnecessária de
        antibióticos, necessidade de encaminhamento ao
        otorrinolaringologista) associados ao uso do modelo,
        limitando-se a métricas de acurácia diagnóstica.
\end{itemize}

% --- 6. CONTRIBUIÇÃO PARA O LIVRO ---
\subsection{Contribuição para o Livro}
Este artigo fundamenta a Seção 8.5 do Capítulo 8, que aborda
especificamente o diagnóstico automatizado de sinusite maxilar com
deep learning. O modelo YOLOv8-MS com precisão de 97,1\% e F1-score
de 95,4\% é utilizado como benchmark primário na discussão sobre
detecção de patologias sinusais em CBCT. A diferenciação entre cisto
de retenção, espessamento mucoso e sinusite odontogênica é central
para a Seção 8.5.3 (Diagnóstico diferencial automatizado), onde o
fluxo de decisão do modelo é comparado com o raciocínio clínico
tradicional. A capacidade de classificar gravidade (AUC = 0,942)
subsidia a subseção sobre estadiamento automatizado e planejamento
cirúrgico. A comparação com radiologistas (modelo $>$ generalista
$\approx$ especialista) é citada na Seção 8.6 como evidência de que
sistemas de IA bem calibrados podem atuar como equalizadores de
qualidade diagnóstica, especialmente em serviços com baixa
disponibilidade de especialistas. A limitação monocêntrica é
discutida no livro como justificativa para a inclusão de validação
multicêntrica obrigatória no protocolo SDD.

% --- 7. ANÁLISE CRÍTICA ---
\subsection{Análise Crítica}
\begin{itemize}
  \item \textbf{Validade interna:} Alta -- Amostra robusta (1.080
        imagens, 2.158 seios), anotação por consenso de especialistas,
        arquitetura comparada com variantes YOLO e outras CNNs,
        análise estatística com IC 95\% e testes de significância.
  \item \textbf{Validade externa:} Baixa a Média -- Validação externa
        limitada a 200 imagens de 2 centros (ambos coreanos, mesmo
        perfil populacional e equipamento similar); ausência de
        validação em populações ocidentais, africanas ou com outras
        marcas de CBCT impede a generalização global.
  \item \textbf{Reprodutibilidade:} Média -- Código do modelo
        disponível no GitHub (repositório público, licença MIT), mas
        os dados brutos (CBCT) não podem ser compartilhados por
        restrições éticas; o pré-processamento e os pesos treinados
        estão disponíveis mediante solicitação.
  \item \textbf{Relevância clínica:} Alta -- A sinusite maxilar é a
        patologia mais frequente dos seios paranasais, com prevalência
        de 25--45\% em candidatos a implantes na região posterior
        maxilar; o diagnóstico precoce e preciso é crucial para o
        planejamento cirúrgico e prevenção de complicações.
  \item \textbf{Conflito de interesses:} Não -- Declaração de ausência
        de conflitos; financiamento por agências públicas (NRF Coreia,
        JSPS Japão, CAPES Brasil).
  \item \textbf{Viés identificado:} Viés de espectro — a amostra
        incluiu predominantemente sinusites moderadas a graves
        (72,8\% dos casos de sinusite), com sub-representação de
        casos leves e sutis que são justamente os de maior desafio
        diagnóstico na prática clínica.
\end{itemize}

% --- 8. CITAÇÕES RELEVANTES ---
\subsection{Citações Relevantes}
\begin{citacao}
``The YOLOv8-MS model achieved an overall precision of 97.1\% for
maxillary sinusitis detection, with class-specific precisions of
96.9\% (normal), 95.2\% (retention cyst), and 99.2\% (sinusitis).
The model outperformed general radiologists (86.4\% accuracy) and
matched expert-level performance (96.8--97.3\%), while reducing
inference time to 28 ms per image — a 350-fold improvement over
manual annotation.'' (p. 458).
\end{citacao}

% --- 9. MAPEAMENTO SDD ---
\subsection{Mapeamento SDD}
\textbf{Problema:} Diagnóstico de sinusite maxilar em CBCT com alta
variabilidade interobservador, necessidade de expertise radiológica
especializada e risco de subdiagnóstico de sinusite odontogênica,
que pode comprometer o planejamento de cirurgias reconstrutoras e
implantes na região posterior maxilar. \\
\textbf{Entrada:} Cortes coronais de CBCT da região maxilar (512
$\times$ 512 pixels, 16-bit, 0,2--0,4 mm voxel); 2.158 seios
maxilares anotados por consenso. \\
\textbf{Saída:} Classificação triádica (normal / cisto de retenção /
sinusite); classificação de gravidade (leve / moderada / grave);
diferenciação etiológica (odontogênica vs. não odontogênica); mapa
de segmentação do seio maxilar (DSC $\geq$ 0,93). \\
\textbf{Critérios:} Precisão $>$ 95\%; F1-score $>$ 0,94; AUC $>$ 0,92;
DSC $>$ 0,92; tempo de inferência $<$ 50 ms/imagem; validação externa
mínima 2 centros. \\
\textbf{Confiança:} 0,78 — Modelo com desempenho excelente e potencial
clínico imediato, mas a validade externa limitada (apenas centros
coreanos) e o viés de espectro reduzem a confiança na generalização
para populações globais e equipamentos heterogêneos. A ausência de
validação prospectiva também contribui para a pontuação moderada.

% --- 10. REFERÊNCIA ABNT ---
\subsection{Referência ABNT}
\begin{verbatim}
KANG, Dongwoo et al. Deep learning for automated detection of maxillary
sinusitis in cone-beam computed tomography images. Oral Surgery, Oral
Medicine, Oral Pathology and Oral Radiology, New York, v. 138, n. 5,
p. 447-462, 2024. DOI: 10.1016/j.oooo.2024.01.018.
\end{verbatim}
\endinput
